\chapter{Implementation and Prototype}
\label{ch:implementation}

The previous chapters describe a design (Chapter~\ref{ch:system-design}), a measurement
protocol (Chapter~\ref{ch:methodology}) and what the measurement found
(Chapter~\ref{ch:results}). This chapter describes the artefact, the HORUS system
(\S\ref{sec:contributions}): the software that runs the pipeline, produces the numbers, and
lets a human look at what the system did.

It is deliberately short on inventory and long on the few structural choices that had
methodological consequences. Three of them recur: configuration is data, not code,
every experiment is resumable and records its own provenance, and every research surface
recomputes what it displays from the stored per-field outcomes instead of showing a summary
written earlier. Each of those choices is the reason some claim elsewhere in this thesis can
be checked instead of taken on trust.

\section{Software Architecture}
\label{sec:sys-architecture}

The implementation is a Python package with a strictly one-directional dependency structure,
which matters less as a matter of taste than as a matter of what can be tested in isolation.

At the base sit the pure functions: value normalisation, ground-truth extraction from the
embedded invoice representation, and the scoring rules. These import nothing from the layers
above them, take data in and return data out, and are consequently testable without a model,
a corpus or a network. Above them sit the pipeline stages --- rasterisation, the vision-model
wrapper, the transcript-to-record structurer, the typed schema with its repair pass. Above
those sits the orchestration harness, and beside everything sits configuration and run
tracking.

One detail of that layering is worth naming because it prevented a specific class of error.
The prediction-side normalisers are imported by both the scorer and the typed schema, and
import neither. Before that arrangement, canonicalisation existed in two places, and the two
could disagree --- which is exactly the failure mode Chapter~\ref{ch:measurement-validity}
documents when a rule was applied to one side of a comparison and not the other. A single
implementation with two consumers makes the asymmetric case impossible to write by accident.

\subsection{Configuration as data}

Every knob of every experiment --- model identifiers, dataset paths, seeds, decoding
settings, prompt strings, token budgets, thresholds, tracking metadata --- lives in a
versioned configuration file. Experiments take exactly one argument: the path to that file.
No experiment reads a constant from its own source.

The configuration is a typed schema validated before anything else happens. Two properties of
that schema do real work. Unknown keys are rejected, not ignored, so a misspelt setting
fails immediately instead of silently taking its default --- a failure mode that is
particularly insidious in machine-learning code, where a run with a wrong setting still
completes and still produces a number. And validation happens before any model is loaded, so a
malformed configuration costs a second instead of the several minutes it takes to read
quantised weights into memory.

Overlay composition is supported: a small file can override selected keys of a base
configuration, deep-merging into it, so a development variant of a sweep does not duplicate the
full specification. Environment variables can override individual settings, which is how
credentials stay out of the repository.

\subsection{The orchestration harness}

The evaluation sweeps in this thesis are cross-products --- several models against several
dozen documents --- executed on a laptop. Three properties follow from that setting, not
from architectural preference.

\textbf{Every cell is an independent unit of work.} Each (model, document) pair is rasterised,
run, scored and recorded on its own. A failure in one cell is caught, recorded with its error
class, and does not stop the sweep; a model that fails to load marks its whole row with a
reason and the harness moves to the next model. A sweep that aborts on the first failure would
have made the cohort survey of \S\ref{sec:results-preliminary} impossible, because several of
the candidate models failed in exactly that way.

\textbf{Sweeps are resumable.} Before running a cell the harness asks the tracking store
whether that cell already completed, and skips it if so. Interrupting a multi-hour sweep and
restarting it therefore costs only the cell in flight. This is not a convenience feature: on
this hardware a full sweep does not reliably fit into one uninterrupted session, and without
resumability the practical choice would have been between smaller sweeps and unreliable ones.

\textbf{Every cell records its own provenance.} Each run carries the configuration that
produced it, the model identifier, the document identifier, the per-field outcomes and the raw
model output as stored artefacts. The tracking metadata lives in a local database and the
artefacts on the local filesystem; nothing leaves the machine. This is what makes the
re-scoring of Chapter~\ref{ch:measurement-validity} possible, and it is also what makes a
number in this thesis traceable to the run that produced it, not to a note.

\subsection{Two training executors, one data path}

The adaptation study of \S\ref{sec:results-adaptation} had to move off the local machine, for
reasons \S\ref{sec:method-repro} states. That created a risk worth naming: two training paths
means two opportunities for the training data to differ, and a difference there would have
made the two arms incomparable and the whole grid uninterpretable.

The implementation avoids that by sharing everything except the executor. Both paths read the
same sealed split and verify its recorded fingerprints, carve the same selection slice, load
the same instruction string from the same configuration file, and construct their training
targets with the same code, including the consistency gate that discards any example whose
constructed target does not itself score perfectly against the reference. Only the code that
performs the optimisation differs. The checkpoint-selection rule is likewise delegated to
standard trainer configuration --- select on the lowest loss over the held-back slice, evaluate
once per epoch --- instead of implemented by hand, so the executed rule cannot drift from the
registered one, and the sealed validation set is never passed to the trainer at all.

\subsection{Verification}
\label{sec:impl-verification}

The project carries a test suite of 1{,}265 tests. Its composition reflects where the risk
actually sits: the largest file by a clear margin covers the scoring rules, and the orchestration
harness, the reference extraction, and the typed schema with its value normalisers and repair
pass follow, because those are the components whose defects corrupt results silently instead of
loudly. Several tests exist specifically to pin down invariants that were violated once --- that
the field registry has the size the reported figures assume, that normalisation is applied to both
sides of a comparison, and that no concrete value from the corpus appears inside a prompt.

Tests that require the corpus, or that require Apple-Silicon inference, skip themselves by
predicate instead of being excluded by a marker or deleted. The same suite therefore runs
locally with the corpus present, locally without it, and in continuous integration, with the
skips visible in the output, not hidden in configuration.

Continuous integration runs on every push to the main branch and every pull request against it:
linting, formatting, type-checking with a cleared cache, and the full suite. The cleared
type-check cache is a response to a specific incident --- a warm cache silently masked an import
error after a configuration change --- and is documented as such, not as a general caution.

\subsection{The operational interface}
\label{sec:impl-assets}

Every operation in the project is a named build target: install, test, lint, type-check,
generate the synthetic corpus and validate it against an independent implementation, run a
smoke test for one model, run a full sweep, browse the tracking store, run the dashboard,
regenerate every table and figure in this manuscript, and build this manuscript.

That is a deliberate choice about where operational knowledge lives. A command a researcher
remembers is a command that stops working when they forget it; a target in the build file is
discoverable, reviewable in a diff, and identical for every invocation. The tables and figures
in Chapters~\ref{ch:results} and~\ref{ch:measurement-validity} are generated by one such
target from committed evaluation output, which is what makes the prohibition on transcribing
numbers into the manuscript enforceable instead of aspirational.

\section{Observability and Demo Application}
\label{sec:sys-app}

A per-field evaluation produces a great deal of structure that a table cannot show: which
document, which field, which value, against which page region, in which model's output. The
prototype application exists to make that structure inspectable. It runs locally, reads from
the same tracking store the sweeps write to, and has the five surfaces of
Figure~\ref{fig:app-surfaces}.

\begin{figure}[htbp]
  \centering
  %% The five surfaces of the prototype application (\S\ref{sec:sys-app}) around the
%% local tracking store: three read-only research surfaces that recompute at display
%% time, one live-extraction surface that runs the pipeline, and the adjudication
%% surface that produced the held-out answer key. The read/write asymmetry is the
%% design content: adjudication writes to a separate location and never overwrites
%% the channels it reads.
%%
%% Hand-authored TikZ (not generated): it depicts the application's design, not a
%% measurement.
%%
%% Input from a chapter with:  %% The five surfaces of the prototype application (\S\ref{sec:sys-app}) around the
%% local tracking store: three read-only research surfaces that recompute at display
%% time, one live-extraction surface that runs the pipeline, and the adjudication
%% surface that produced the held-out answer key. The read/write asymmetry is the
%% design content: adjudication writes to a separate location and never overwrites
%% the channels it reads.
%%
%% Hand-authored TikZ (not generated): it depicts the application's design, not a
%% measurement.
%%
%% Input from a chapter with:  %% The five surfaces of the prototype application (\S\ref{sec:sys-app}) around the
%% local tracking store: three read-only research surfaces that recompute at display
%% time, one live-extraction surface that runs the pipeline, and the adjudication
%% surface that produced the held-out answer key. The read/write asymmetry is the
%% design content: adjudication writes to a separate location and never overwrites
%% the channels it reads.
%%
%% Hand-authored TikZ (not generated): it depicts the application's design, not a
%% measurement.
%%
%% Input from a chapter with:  \input{figures/app-surfaces.tex}

\begin{tikzpicture}[
    font=\small,
    store/.style={
      draw=horusaccent, line width=1.0pt, rounded corners=2pt,
      fill=horuswash, text width=3.2cm, align=center, inner sep=5pt,
    },
    surface/.style={
      draw=horusaccent, line width=0.7pt, rounded corners=2pt,
      fill=white, text width=3.2cm, align=center,
      font=\scriptsize, inner sep=4pt,
    },
    writer/.style={
      draw=horuswarn, line width=0.7pt, rounded corners=2pt,
      fill=white, text width=3.2cm, align=center,
      font=\scriptsize, inner sep=4pt,
    },
    readflow/.style={-{Stealth[length=2.0mm]}, line width=0.6pt, draw=horusaccent},
    writeflow/.style={-{Stealth[length=2.0mm]}, line width=0.7pt, draw=horuswarn,
      dashed},
    grp/.style={font=\scriptsize\bfseries, text=horusink, align=center},
  ]

  \node[store] (store) {\textbf{local tracking store}\\[1pt]\scriptsize per-cell
    provenance: configuration, page image, transcript, raw output, parsed record,
    per-field outcomes};

  %% Three read-only surfaces on the left; the shared design rule sits under them.
  \node[surface, above left=6mm and 10mm of store] (overview)
    {\textbf{overview}\\headline accuracy per approach};
  \node[surface, left=10mm of store] (explorer)
    {\textbf{invoice explorer}\\page image, transcript, record, per-field verdict};
  \node[surface, below left=6mm and 10mm of store] (compare)
    {\textbf{approach comparison}\\methods side by side, per-field failure breakdown};

  \draw[readflow] (store.west) ++(-1mm,0) -- (explorer.east);
  \draw[readflow] (store.north west) -- (overview.east);
  \draw[readflow] (store.south west) -- (compare.east);

  \node[grp, above=2mm of overview, text width=6.8cm, align=left, anchor=south west,
        xshift=-2mm] {read-only: recompute from stored outcomes at display time ---
    a surface that recomputes cannot show a stale number};

  %% Live extraction on the upper right: same configuration, same resolution, never
  %% scores. Adjudication on the lower right: the one writer, quarantined.
  \node[surface, above right=6mm and 10mm of store] (live)
    {\textbf{live extraction}\\runs the pipeline on this machine, at the evaluated
    resolution; never scores --- an upload has no reference};
  \draw[readflow] (live.south west) -- (store.north east);

  \node[writer, below right=6mm and 10mm of store] (adjud)
    {\textbf{adjudication surface}\\248 escalated cells, ranked by damage; sign-off
    writes provenance to a separate location, never over the channels it reads};
  \draw[readflow] (store.south east) -- (adjud.north west);
  \draw[writeflow] (adjud.south) -- ++(0,-4mm) -| node[pos=0.25, below,
    font=\scriptsize, text=horuswarn] {signed-off answer key (separate location)}
    ++(-2.0cm,0);

\end{tikzpicture}


\begin{tikzpicture}[
    font=\small,
    store/.style={
      draw=horusaccent, line width=1.0pt, rounded corners=2pt,
      fill=horuswash, text width=3.2cm, align=center, inner sep=5pt,
    },
    surface/.style={
      draw=horusaccent, line width=0.7pt, rounded corners=2pt,
      fill=white, text width=3.2cm, align=center,
      font=\scriptsize, inner sep=4pt,
    },
    writer/.style={
      draw=horuswarn, line width=0.7pt, rounded corners=2pt,
      fill=white, text width=3.2cm, align=center,
      font=\scriptsize, inner sep=4pt,
    },
    readflow/.style={-{Stealth[length=2.0mm]}, line width=0.6pt, draw=horusaccent},
    writeflow/.style={-{Stealth[length=2.0mm]}, line width=0.7pt, draw=horuswarn,
      dashed},
    grp/.style={font=\scriptsize\bfseries, text=horusink, align=center},
  ]

  \node[store] (store) {\textbf{local tracking store}\\[1pt]\scriptsize per-cell
    provenance: configuration, page image, transcript, raw output, parsed record,
    per-field outcomes};

  %% Three read-only surfaces on the left; the shared design rule sits under them.
  \node[surface, above left=6mm and 10mm of store] (overview)
    {\textbf{overview}\\headline accuracy per approach};
  \node[surface, left=10mm of store] (explorer)
    {\textbf{invoice explorer}\\page image, transcript, record, per-field verdict};
  \node[surface, below left=6mm and 10mm of store] (compare)
    {\textbf{approach comparison}\\methods side by side, per-field failure breakdown};

  \draw[readflow] (store.west) ++(-1mm,0) -- (explorer.east);
  \draw[readflow] (store.north west) -- (overview.east);
  \draw[readflow] (store.south west) -- (compare.east);

  \node[grp, above=2mm of overview, text width=6.8cm, align=left, anchor=south west,
        xshift=-2mm] {read-only: recompute from stored outcomes at display time ---
    a surface that recomputes cannot show a stale number};

  %% Live extraction on the upper right: same configuration, same resolution, never
  %% scores. Adjudication on the lower right: the one writer, quarantined.
  \node[surface, above right=6mm and 10mm of store] (live)
    {\textbf{live extraction}\\runs the pipeline on this machine, at the evaluated
    resolution; never scores --- an upload has no reference};
  \draw[readflow] (live.south west) -- (store.north east);

  \node[writer, below right=6mm and 10mm of store] (adjud)
    {\textbf{adjudication surface}\\248 escalated cells, ranked by damage; sign-off
    writes provenance to a separate location, never over the channels it reads};
  \draw[readflow] (store.south east) -- (adjud.north west);
  \draw[writeflow] (adjud.south) -- ++(0,-4mm) -| node[pos=0.25, below,
    font=\scriptsize, text=horuswarn] {signed-off answer key (separate location)}
    ++(-2.0cm,0);

\end{tikzpicture}


\begin{tikzpicture}[
    font=\small,
    store/.style={
      draw=horusaccent, line width=1.0pt, rounded corners=2pt,
      fill=horuswash, text width=3.2cm, align=center, inner sep=5pt,
    },
    surface/.style={
      draw=horusaccent, line width=0.7pt, rounded corners=2pt,
      fill=white, text width=3.2cm, align=center,
      font=\scriptsize, inner sep=4pt,
    },
    writer/.style={
      draw=horuswarn, line width=0.7pt, rounded corners=2pt,
      fill=white, text width=3.2cm, align=center,
      font=\scriptsize, inner sep=4pt,
    },
    readflow/.style={-{Stealth[length=2.0mm]}, line width=0.6pt, draw=horusaccent},
    writeflow/.style={-{Stealth[length=2.0mm]}, line width=0.7pt, draw=horuswarn,
      dashed},
    grp/.style={font=\scriptsize\bfseries, text=horusink, align=center},
  ]

  \node[store] (store) {\textbf{local tracking store}\\[1pt]\scriptsize per-cell
    provenance: configuration, page image, transcript, raw output, parsed record,
    per-field outcomes};

  %% Three read-only surfaces on the left; the shared design rule sits under them.
  \node[surface, above left=6mm and 10mm of store] (overview)
    {\textbf{overview}\\headline accuracy per approach};
  \node[surface, left=10mm of store] (explorer)
    {\textbf{invoice explorer}\\page image, transcript, record, per-field verdict};
  \node[surface, below left=6mm and 10mm of store] (compare)
    {\textbf{approach comparison}\\methods side by side, per-field failure breakdown};

  \draw[readflow] (store.west) ++(-1mm,0) -- (explorer.east);
  \draw[readflow] (store.north west) -- (overview.east);
  \draw[readflow] (store.south west) -- (compare.east);

  \node[grp, above=2mm of overview, text width=6.8cm, align=left, anchor=south west,
        xshift=-2mm] {read-only: recompute from stored outcomes at display time ---
    a surface that recomputes cannot show a stale number};

  %% Live extraction on the upper right: same configuration, same resolution, never
  %% scores. Adjudication on the lower right: the one writer, quarantined.
  \node[surface, above right=6mm and 10mm of store] (live)
    {\textbf{live extraction}\\runs the pipeline on this machine, at the evaluated
    resolution; never scores --- an upload has no reference};
  \draw[readflow] (live.south west) -- (store.north east);

  \node[writer, below right=6mm and 10mm of store] (adjud)
    {\textbf{adjudication surface}\\248 escalated cells, ranked by damage; sign-off
    writes provenance to a separate location, never over the channels it reads};
  \draw[readflow] (store.south east) -- (adjud.north west);
  \draw[writeflow] (adjud.south) -- ++(0,-4mm) -| node[pos=0.25, below,
    font=\scriptsize, text=horuswarn] {signed-off answer key (separate location)}
    ++(-2.0cm,0);

\end{tikzpicture}

  \caption[The prototype's five surfaces]{The prototype's five surfaces around the local
  tracking store. The three research surfaces recompute from stored per-field outcomes at
  display time; the live-extraction surface runs the measured pipeline; the adjudication
  surface is the only writer, and it writes to a separate location so the inputs to
  adjudication stay independent of its outputs.}
  \label{fig:app-surfaces}
\end{figure}

\textbf{Three read-only research surfaces.} An overview reports headline accuracy per
approach; an invoice explorer shows one document for one approach --- the page image the model
saw, the raw text it emitted, the parsed record, and a colour-coded verdict for every field
against the reference; an approach comparison places the methods side by side on the documents
they share, with a per-field breakdown of where each one fails.

One design rule governs all three: they recompute their figures from the stored per-field
outcomes at display time instead of reading a summary number that was written earlier. The
reason is the failure mode of Chapter~\ref{ch:measurement-validity} --- a stored summary
computed under a superseded scoring rule looks exactly like a current one. A surface that
recomputes cannot show a stale number, and during the corrections it repeatedly served as the
first place a discrepancy became visible.

\textbf{A live extraction surface.} One page departs from read-only: upload a document, choose
a method, and the pipeline runs on this machine. Two properties make it honest rather than a
demonstration trick. It rasterises at the same resolution the evaluation uses, so what it shows
is what was measured rather than a favourable rendering. And it never scores, because an
uploaded document has no reference --- it presents the extraction for human inspection and
says so. Model identifiers, prompts and token budgets are read from the same configuration
files the experiments use, so the demonstration cannot drift from the measured system.

\textbf{An adjudication surface.} The fifth page is the annotation interface behind the
held-out answer key of \S\ref{sec:method-heldout-gt}, and it is the clearest case in this
project of an interface decision determining whether a methodological step was feasible at all.

The naive form of that interface is a form with every field for every document: 34 fields
across 39 documents is 1{,}326 decisions. That review does not get finished, and an unfinished
review produces an answer key whose warrant is unclear --- which is precisely how an earlier
reported figure came to rest on an unverified key and had to be withdrawn. The implemented
surface instead shows only the cells the automatic combiner could not settle, ranked so that
the cells where an undetected error would do the most damage come first, and presents for each
one every channel's reading together with the page text that supports it. That is 248
decisions, and 248 decisions can be finished.

Two safeguards are built into it. Sign-off writes to a separate location carrying a provenance
record per cell, and never overwrites the files of the channels the combiner reads --- so the
inputs to adjudication remain independent of its outputs. And the page runs no model: every
reading it displays was produced earlier and is read from disk, so the act of reviewing cannot
introduce a new source of evidence.

\section{Deployment}
\label{sec:sys-deploy}

This section states what exists, because the honest answer is narrower than an implementation
chapter usually reports.

\textbf{What exists} is a reproducible local development and evaluation environment: a locked
dependency manifest resolved by a single command, a pinned interpreter version, the build
targets described above, and continuous integration that verifies the whole of it on every
change. On the floor machine --- one laptop, no discrete accelerator, no network dependency on
the inference path --- a checkout runs the delivered pipeline and every scoring and re-scoring
pass reported in this thesis; the training and native-resolution reading passes require the
target class (\S\ref{sec:method-repro}). One external
dependency sits outside the Python environment: an independent validator used to cross-check
the reference extraction is fetched on demand and verified against a pinned cryptographic
hash, so the cross-check cannot silently become a different tool.

\textbf{What does not exist} is packaging for anyone other than a developer. There is no
container image, no service interface, no installer, and no operational documentation for a
firm's information-technology department. The system is operated from a source checkout by
someone who can read a build file.

That is a scope decision, not an omission, and the reasoning is the same as in
\S\ref{sec:design-scope}. Packaging serves adoption; nothing in the research questions of
\S\ref{sec:research-questions} depends on it, and the effort it would have consumed went into
the measurement apparatus instead --- which, given what Chapter~\ref{ch:measurement-validity}
found, was the correct allocation. What a deployable version would additionally require is
listed in \S\ref{sec:fw-production}, and the prototype's review surface is the part of it that
already exists, because the evaluation needed it first.
